\uuid{9uC4}
\titre{Nombres premiers}
\auteur{bodin ; croizat ; sacre}
\organisation{exo7}
\licence{CC BY-NC-SA 4.0}
\datecreate{2022-03-08}
\niveau{L1}
\module{Algèbre}
\chapitre{Arithmétique dans Z}
\sousChapitre{Nombres premiers, nombres premiers entre eux}
\difficulte{2}
\theme{Nombres premiers}
\exo7qcmid{hilisit.arithmetique.032}

\contenu{
\begin{question}
\qclassification{103.04}
    Soit $n$ un entier tel que $n^2-1$ est un multiple de $11$. Quelles sont les affirmations vraies ?
    \begin{answers}
        \bad{$11$ divise $n-1$.}
        \bad{$11$ divise $n+1$.}
        \good{($11$ divise $n-1$) ou ($11$ divise $n+1$).}
        \bad{($11$ divise $n-1$) et ($11$ divise $n+1$).}
    \end{answers}
    \begin{explanations}
    $11$ est premier et divise $n^2-1 = (n-1)(n+1)$. D'après le lemme d'Euclide, soit $11 | (n-1)$, soit $11 | (n+1)$.\\
    Pour $n=10$, on a $11 | (10^2-1)$ mais $11$ ne divise pas $10-1=9$.\\
    Pour $n=12$, on a $11 | 12^2 - 1$ mais $11$ ne divise pas $12+1=13$.\\
    Et si $11$ divisait $n-1$ et $n+1$ alors $11$ diviserait  $n+1-(n-1) = 2$.
    \end{explanations}
\end{question}
}
